\documentclass[12pt,a4paper]{article}

\usepackage{fontspec}
\usepackage{xeCJK}
\setCJKmainfont{Hiragino Mincho ProN}
\setCJKsansfont{Hiragino Kaku Gothic ProN}
\setCJKmonofont{Hiragino Kaku Gothic ProN}
\usepackage[margin=2.5cm]{geometry}
\usepackage{amsmath,amssymb}
\usepackage{hyperref}
\usepackage{booktabs}
\usepackage{tcolorbox}
\usepackage{enumitem}

\tcbuselibrary{breakable,skins}

\newtcolorbox{keybox}[1][]{
  colback=black!3, colframe=black!50,
  fonttitle=\bfseries, title={#1},
  breakable, sharp corners
}

\newtcolorbox{questionbox}{
  colback=blue!3, colframe=blue!40,
  fonttitle=\bfseries, title={素朴な疑問},
  breakable, sharp corners
}

\newtcolorbox{analogybox}{
  colback=green!3, colframe=green!40,
  fonttitle=\bfseries, title={たとえ話},
  breakable, sharp corners
}

\newtcolorbox{problembox}{
  colback=red!3, colframe=red!40,
  fonttitle=\bfseries, title={ここが問題},
  breakable, sharp corners
}

\setlength{\parskip}{0.8em}
\setlength{\parindent}{0pt}

\begin{document}

\begin{center}
{\LARGE\bfseries 高校生のための入門ガイド}\\[0.5cm]
{\Large\bfseries ワープドライブと数学のあいだ}\\[0.3cm]
{\large --- グロタンディークの「層」で物理法則をハックする ---}\\[1.5cm]
{\large Wright Brothers}\\[0.3cm]
{2026年4月}
\end{center}

\vspace{1cm}

\tableofcontents

\newpage

% ============================================================================
\section{この本が何の話をしているか}
% ============================================================================

SF映画に出てくる「ワープドライブ」。
光より速く宇宙を移動する夢の技術。
実は、アインシュタインの一般相対論の中に、
ワープを実現する数学的な解（アルクビエレ計量）が1994年に見つかっている。

ただし、1つだけ「絶対に必要な材料」がある。

それは\textbf{「負のエネルギー」}。

この負のエネルギーを作ることは、物理学のルール（弱エネルギー条件）で
基本的に禁止されている。
---ように見える。

本当に「絶対に」禁止されているのか？
それとも、ルールの書き方を変えたら、禁止が消えるのか？

この論文は、「ルールの書き方を変える」話だ。
そして、そのためにグロタンディークという数学者が
60年前に発明した「層」（sheaf）という道具を使う。

% ============================================================================
\section{ワープドライブの材料：負のエネルギー}
% ============================================================================

\subsection{普通のエネルギーは「正」}

日常のすべてのエネルギーは正の値を持つ。
電気も、熱も、光も、質量（$E=mc^2$）も、全部正。

物理学には「弱エネルギー条件（WEC）」というルールがある：

\begin{keybox}[弱エネルギー条件 (WEC)]
どんな観測者がどんな場所で測っても、
エネルギー密度は $\rho \geq 0$（ゼロ以上）でなければならない。
\end{keybox}

このルールは「当たり前」に見える。
エネルギーがマイナスって、何？

\subsection{カシミール効果：真空のエネルギーは「負」になれる}

1948年、カシミールという物理学者が
驚くべきことを理論的に予測した（のちに実験で確認）：

2枚の金属板をごく近く（ナノメートル）に置くと、
板の間の「真空」のエネルギーが\textbf{負になる}。
板が互いに引き合う力が生まれる。

\begin{analogybox}
海を想像してほしい。
波は無限にいろんな波長で立っている。
でも、狭い入り江の中では、入り江に「はまる」波長の波しか立てない。
入り江の外より波が少ない = エネルギーが少ない = 外を基準にすると「負」。
\end{analogybox}

つまり、「境界条件を変える」と、真空のエネルギーは変わる。
負にすらなれる。

でも、カシミール効果で得られる負のエネルギーは
あまりにも小さすぎて、ワープには全然足りない。
$10^{69}$倍（1のうしろにゼロが69個！）足りない。

\begin{problembox}
カシミール効果で負のエネルギーは作れるが、
ワープに必要な量には $10^{69}$ 桁足りない。

これが「スケーリング問題」。
\end{problembox}

% ============================================================================
\section{真空エネルギーの計算とリーマンゼータ関数}
% ============================================================================

\subsection{無限の足し算をどう「正しく」やるか}

真空のエネルギーを計算すると、こういう式が出てくる：

$$
E = \frac{\hbar\omega_0}{2}(1 + 2 + 3 + 4 + 5 + \cdots)
$$

この無限和は発散する（無限大になる）。
でも物理学者は、「ゼータ関数正則化」という技を使って、
この和に有限の値を割り当てる。

$$
1 + 2 + 3 + 4 + \cdots \longrightarrow \zeta(-1) = -\frac{1}{12}
$$

これは「$1+2+3+\cdots = -1/12$」という有名な話。
インチキに聞こえるかもしれないが、
この値を使って計算したカシミール力は
実験で測定でき、ぴったり一致する。

つまり、真空のエネルギーは
\textbf{リーマンゼータ関数} $\zeta(s)$ の
特殊な値で決まる。

\subsection{オイラー積：素数の寄与}

リーマンゼータ関数には、素数を使った「別の書き方」がある。
これをオイラー積という：

$$
\zeta(s) = \frac{1}{1-2^{-s}} \times \frac{1}{1-3^{-s}} \times \frac{1}{1-5^{-s}} \times \frac{1}{1-7^{-s}} \times \cdots
$$

\begin{keybox}[オイラー積の意味]
リーマンゼータ関数は、\textbf{各素数がそれぞれ独立に}寄与する
「チャンネル」の積として分解できる。

素数2のチャンネル $\times$ 素数3のチャンネル $\times$ 素数5のチャンネル $\times \cdots$
\end{keybox}

これは驚くべきことだ。
真空のエネルギーが、素数という「数学の構造」で決まっている。

% ============================================================================
\section{「素数チャンネルのミュート」で符号が反転する}
% ============================================================================

では、素数2のチャンネルを「ミュート」（消音）したらどうなるか？

オイラー積から素数2の因子だけを外す：

$$
\zeta_{\neg 2}(s) = \zeta(s) \times (1 - 2^{-s})
$$

3次元の真空エネルギーに対応する $s = -3$ で計算すると：

$$
\zeta_{\neg 2}(-3) = \frac{1}{120} \times (1 - 2^3) = \frac{1}{120} \times (-7) = -\frac{7}{120}
$$

\begin{keybox}[符号反転定理]
もとの真空エネルギー $\zeta(-3) = +1/120$（正）。

素数2をミュートした後の真空エネルギー $\zeta_{\neg 2}(-3) = -7/120$（\textbf{負！}）。

符号が $+$ から $-$ に反転した。
\end{keybox}

しかも、これは素数2に限らない。
\textbf{どの素数をミュートしても、必ず負になる。}

$$
\zeta_{\neg p}(-3) = \frac{1-p^3}{120} < 0 \quad (\text{なぜなら } p^3 \geq 8 > 1)
$$

負の真空エネルギー = WEC違反 = ワープの材料。

\begin{questionbox}
ちょっと待って。「素数をミュートする」って、物理的に何をするの？

--- まさにそこが問題。次の章で説明する。
でもその前に、もっと大きな問題がある。
\end{questionbox}

% ============================================================================
\section{10の69乗倍問題：カントールの壁}
\label{sec:scaling}
% ============================================================================

素数をミュートすれば負のエネルギーが作れる。
しかし、実験室で作れる負のエネルギーは
$\sim 10^{-22}$ ジュール（原子1個分のエネルギーよりはるかに小さい）。

ワープドライブに必要なのは $\sim 10^{47}$ ジュール。

差は $10^{69}$ 倍。宇宙に存在する原子の数（$10^{80}$）に近い桁数。

\begin{problembox}
小さな実験で負のエネルギーが作れても、
ワープに必要な量に到達するには $10^{69}$ 倍の
「スケールアップ」が必要。

これが解決不可能に見える。

---ここまでが、「ふつうの」物理学と数学の考え方で
たどり着ける限界。
\end{problembox}

ここで、この論文が問いかける：

\textbf{この「$10^{69}$倍足りない」という問題は、
本当に「量」の問題なのか？}

もし「量」の問題ではなく「構造」の問題だとしたら？

つまり：
\begin{itemize}
\item 旧い考え方：負のエネルギーを $10^{69}$ 倍に増幅しなければならない
\item 新しい考え方：WECという「ルール」自体を無効化すればよい
\end{itemize}

ルールの無効化は、量の蓄積とは全く異なる種類の操作だ。
電気のスイッチをオフにするのに、力は要らない。

この「新しい考え方」を可能にするのが、
グロタンディークの数学だ。

% ============================================================================
\section{カントール vs グロタンディーク：数学の2つの世界観}
\label{sec:cantor_vs_grothendieck}
% ============================================================================

\subsection{カントールの世界観：点の集合}

19世紀のカントール以来、数学は「集合論」を基礎にしている。
空間 = 点の集まり。

この見方で真空エネルギーを計算すると：
$$
\zeta(s) = 1^{-s} + 2^{-s} + 3^{-s} + 4^{-s} + \cdots
$$
各自然数を「点」として、その値を足し合わせる。

「素数2のミュート」は、この足し算から
$2^{-s} + 4^{-s} + 6^{-s} + \cdots$（2の倍数の項）を引く。

\begin{analogybox}
バケツの中の水（= エネルギー）から、
コップ1杯分（= 素数2の寄与）を汲み出す。
バケツの中の水の量は「ちょっと減った」。

量が変わるだけ。バケツは同じバケツ。
\end{analogybox}

この考え方では、エネルギーの「量」が問題になる。
少ししか変わらないなら、ワープには程遠い。
→ スケーリング問題（$10^{69}$倍）が発生。

\subsection{グロタンディークの世界観：層の圏}

20世紀後半、グロタンディーク（Alexander Grothendieck）は
数学を根底から書き換えた。

彼の革命的な考え方：

\begin{keybox}[グロタンディークの洞察]
\textbf{空間の本質は「点の集合」ではない。}

空間の本質は、その上に載っている「構造」（層 = sheaf）にある。
同じ点の集合でも、載っている層を変えれば、
全く別の空間になる。
\end{keybox}

\begin{analogybox}
カントール的：東京タワーは鉄の分子の集合。
分子を1つ取り除いても、少し軽くなるだけ。

グロタンディーク的：東京タワーは「設計図」（= 層）で決まっている。
設計図を書き換えれば、同じ鉄で全く別の建物が建つ。
分子の数は関係ない。設計図（構造）が全てを決める。
\end{analogybox}

この見方で真空を見ると：

$$
\zeta(s) = \Spec(\mathbb{Z}) \text{ のゼータ関数}
$$

$\Spec(\mathbb{Z})$（「スペック・ゼット」と読む）は、
全ての素数を「点」として持つ幾何学的空間（スキーム）。
ゼータ関数は、この空間の「形」を記述する不変量。

「素数2のミュート」は、このスキームの\textbf{局所化}：
$$
\Spec(\mathbb{Z}) \longrightarrow \Spec(\mathbb{Z}[1/2])
$$

\begin{keybox}[局所化の意味]
局所化とは、\textbf{空間から素数2の点を取り除く}こと。

点を「集合から引く」のではなく、
\textbf{空間自体を別の空間に変える}。

$\mathbb{Z}$（整数環）を $\mathbb{Z}[1/2]$（2で割れる世界）に変える。

これは「水をコップ1杯汲み出す」のではなく、
「バケツの設計図を書き換えて、底の形を変える」ようなもの。
\end{keybox}

% ============================================================================
\section{なぜグロタンディークでスケーリング問題が消えるのか}
\label{sec:why_grothendieck}
% ============================================================================

ここが論文Bの核心。

\subsection{カントール的計算：値が変わる（量的）}

素数2をミュートすると：
$$
\zeta(-3) = +\frac{1}{120} \quad\longrightarrow\quad \zeta_{\neg 2}(-3) = -\frac{7}{120}
$$

数値が $+1/120$ から $-7/120$ に変わった。
「値」の変化。大きさは $|{-7/120 - 1/120}| = 8/120$。
ワープに必要な量に対して圧倒的に小さい。

→ スケーリング問題は残る。

\subsection{グロタンディーク的解釈：向きが反転する（質的）}

同じ計算を、グロタンディークの目で見る。

$\zeta(-3)$ の値は、K理論（代数的トポロジーの一種）の中で
「ボレル・レギュレータ」という写像の像として解釈できる。

$$
\text{reg}: K_2(\mathbb{Z}) \to \mathbb{R}
$$

この写像の「符号」は、数学的に「格子の向き（orientation）」と呼ばれる
トポロジカルな不変量だ。

\begin{keybox}[レギュレータの向き]
$\zeta(-3) > 0$ → レギュレータの向きが「表」 → WECが成立

$\zeta_{\neg 2}(-3) < 0$ → レギュレータの向きが「裏」 → WECが不成立

向きは「表」か「裏」の二択。大きさは関係ない。
\end{keybox}

\begin{analogybox}
コインの裏表を考えてほしい。

カントール的に見ると：
コインの高さ（表面からテーブルまでの距離）が変わった。
$+1\text{mm}$ から $-1\text{mm}$ に。
「もっと下げるにはどれだけ力が要るか？」→ スケーリング問題。

グロタンディーク的に見ると：
コインがひっくり返った。
表か裏かは「量」ではない。コインを裏返す操作は
1回の離散的な操作であり、コインの重さとは無関係。

「ワープに必要なエネルギーをもっと集めなければ」ではなく、
「コインを裏返す方法を見つければよい」。
\end{analogybox}

\begin{keybox}[論文Bの核心的主張（予想）]
弱エネルギー条件（WEC）は、エネルギーの「量」に関する条件ではなく、
$\Spec(\mathbb{Z})$ 上のレギュレータの「向き」に関する
離散的なトポロジカル不変量である。

局所化 $\mathbb{Z} \to \mathbb{Z}[1/p]$ はこの向きを反転させる。

もしこの予想が正しければ、スケーリング問題は存在しない。
「エネルギーを $10^{69}$ 倍集める」必要はなく、
「向きを1回反転させる」だけでよい。
\end{keybox}

\textbf{注意}：これは「予想」であり、まだ証明されていない。
しかし、数学的に整合的であり、反例は見つかっていない。

% ============================================================================
\section{「向きを反転させる」とは物理的に何をすることか}
% ============================================================================

局所化 $\mathbb{Z} \to \mathbb{Z}[1/p]$ の物理的意味は
5通りの言い方ができる（全て数学的に同じ操作）：

\begin{enumerate}[leftmargin=2em]
\item \textbf{ゲージ化}：
  モード $|n\rangle$ と $|2n\rangle$ を「同じもの」として扱う。
  物理学で言う「ゲージ対称性」の導入。

\item \textbf{余剰次元の解放}：
  弦理論的に言えば、各素数は小さく巻き上がった「余剰次元」。
  素数2の次元を「ほどく」（非コンパクト化する）。

\item \textbf{相転移}：
  物質の相転移（氷→水）のように、
  素数2に対応する「秩序パラメータ」を凍結させる。

\item \textbf{論理の変更}：
  空間の「論理体系」（トポス）を変える。
  WECが「真」だった世界から、WECが「偽」の世界に移動する。
  （これは非常に重要。次の章で詳しく説明する。）

\item \textbf{チャンネルオフ}：
  ラジオの周波数チャンネルを切るように、
  真空の「素数2チャンネル」をオフにする。
\end{enumerate}

どの説明も「エネルギーを集める」話ではない。
「構造を変える」話だ。

実験的に最も現実的なのは、
超伝導量子ビットでの離散ゲージ理論シミュレーション（解釈1）や、
トポロジカル材料での相転移の誘起（解釈3）。

% ============================================================================
\section{トポス量子論：「観測の仕方」が現実を決める}
\label{sec:topos}
% ============================================================================

5つの解釈の中で、4番目の「論理の変更」は
最も奇妙に聞こえるが、実は最も本質的かもしれない。
ここを深掘りする。

\subsection{量子力学の「文脈依存性」}

量子力学には、古典物理学にない根本的な性質がある：

\begin{keybox}[コッヘン-シュペッカーの定理 (1967)]
量子系の物理量は、「どの物理量と一緒に測るか」（= 文脈）に
依存して値が変わる。

つまり：\textbf{測定の「文脈」が現実を決める。}

古典物理学では、りんごの質量は「何と一緒に測るか」によらない。
量子力学では、電子のスピンは「どの軸と一緒に測るか」で変わる。
\end{keybox}

これは「観測者が現実を変える」というオカルトの話ではない。
数学的に厳密に証明された定理だ。

\subsection{トポスとは何か}

グロタンディークが発明した「トポス」は、
この文脈依存性を数学的に扱う枠組みだ。

\begin{analogybox}
RPGゲームを想像してほしい。

\textbf{普通の物理学}：
ゲームの世界に「絶対的な地図」がある。
どのプレイヤーが見ても、山は山、川は川。
命題「この場所に宝箱がある」は、真か偽かの二択。

\textbf{トポス的物理学}：
ゲームの世界は「プレイヤーの視点」ごとに異なる。
戦士が見ると壁がある場所に、魔法使いが見ると通路がある。
命題「この場所に宝箱がある」の真偽は、
\textbf{誰が（= どの文脈で）見るかに依存する}。

しかもどちらも「正しい」。矛盾ではない。
真偽が「真/偽」の2値ではなく、より豊かな構造（ハイティング代数）を持つ。
\end{analogybox}

物理学者のDöring（デーリング）と Isham（イシャム）は2008年に、
量子力学全体をトポスの言語で書き直した。
彼らの理論では：

\begin{itemize}
\item 「文脈」= 同時に測定可能な物理量の組
\item 各文脈ごとに「古典的な世界」が1つ定まる
\item 全ての文脈を集めた圏（= 文脈の圏）上に、層が載る
\item 命題の真偽は、文脈ごとに異なる → ハイティング代数
\end{itemize}

\subsection{WECは「文脈」に依存する}

ここからがワープドライブとの接続。

WEC（弱エネルギー条件）は命題だ：「エネルギー密度 $\rho \geq 0$」。

\begin{problembox}
普通の物理学では、WECは「真」か「偽」かの二択。
そして通常の真空では「真」。終了。ワープ不可能。
\end{problembox}

しかし、トポス的に考えると：

\begin{keybox}[トポス的WEC]
WECの真偽は\textbf{文脈に依存する}。

$\Spec(\mathbb{Z})$ というスキーム上のトポスでは、
WECは「真」と判定される。

$\Spec(\mathbb{Z}[1/p])$ というスキーム上のトポスでは、
同じ命題WECが「偽」と判定されうる。

なぜなら、2つのスキームは異なるトポスを定義し、
異なる内部論理（ハイティング代数）を持つから。
\end{keybox}

\begin{analogybox}
法律のたとえ。

日本では「車は左側通行」は「真」。
アメリカでは「車は左側通行」は「偽」。

同じ命題が、法体系（= 論理体系 = トポス）によって異なる真偽値を持つ。

「左側通行を右側通行に変える」のに、
物理的にすべての道路を作り直す必要はない。
法律（= 論理）を書き換えればよい。

WECも同じ。
$\Spec(\mathbb{Z})$ の「法律」では WEC = 真。
$\Spec(\mathbb{Z}[1/p])$ の「法律」では WEC = 偽。

局所化（$\mathbb{Z} \to \mathbb{Z}[1/p]$）は
「法律の書き換え」= 「トポスの変更」= 「論理のリコンパイル」。
\end{analogybox}

\subsection{なぜこれが「鍵」なのか}

5つの解釈を思い出してほしい：
\begin{enumerate}
\item ゲージ化 → モードの同一視（物理的操作）
\item 余剰次元 → 空間の変形（幾何学的操作）
\item 相転移 → 秩序パラメータの変化（熱力学的操作）
\item \textbf{論理の変更 → トポスの切り替え（論理的操作）}
\item チャンネルオフ → 成分の除去（代数的操作）
\end{enumerate}

1, 2, 3, 5 は全て「何かを物理的にする」話。
4 だけが「何も物理的にしない」話。

トポスの切り替えは、\textbf{物質やエネルギーの操作ではない}。
それは\textbf{「どの文脈で物理法則を評価するか」の変更}だ。

\begin{keybox}[これが「ソースコード書き換え」の本質]
時空の「ソースコード」とは何か？
それは物理法則の集合 = トポスの内部論理。

局所化 $\mathbb{Z} \to \mathbb{Z}[1/p]$ は、
ソースコード（トポス）を別バージョンにリコンパイルすること。

リコンパイル前：WEC = TRUE（コンパイル結果に埋め込み）
リコンパイル後：WEC = FALSE（新しいコンパイル結果）

コンパイルにエネルギーは要らない。
必要なのは「コンパイラに渡すパラメータ」の変更だけ。

→ スケーリング問題が消える理由の最も深い説明。
\end{keybox}

これが、Döring-Isham のトポス量子論と、
グロタンディークの算術幾何学が合流する地点であり、
「観測の文脈が現実を決める」という量子力学の根本原理が、
ワープドライブの鍵になっている理由だ。

% ============================================================================
\section{じゃあ具体的にどうやって素数をミュートするの？}
% ============================================================================

ここまで「素数をミュートする」と言ってきたが、
物理的には何をすればいいのか？
実は10個もの方法が見つかっている。

\subsection{方法A：いらない周波数を吸収する（スペクトルフィルタリング）}

最も直感的な方法。共振器（音や電磁波が反射を繰り返す箱）の中で、
$p$の倍数にあたるモードだけを吸い取る。

\begin{analogybox}
ギターの弦を想像してほしい。弦は基本振動（1倍音）、2倍音、3倍音、4倍音…を同時に出している。

ここに「2の倍数だけ吸収するフィルタ」をかけると、2倍音、4倍音、6倍音が消えて、1倍音、3倍音、5倍音だけが残る。

これが「素数$p=2$のミュート」。
\end{analogybox}

具体的な装置候補：
\begin{itemize}
\item \textbf{超伝導回路 + SQUID}：マイクロ波の共振器にSQUID（超伝導量子干渉素子）を入れて、偶数倍音だけを吸収する。論文Aの実験提案がこれ。コスト：約500万〜1000万円。
\item \textbf{光学周波数コム}：レーザーが等間隔の周波数を出す装置（ノーベル賞技術！）にFabry-P\'erot共振器を組み合わせ、素数の倍数だけ反射させる。コスト：約100万〜500万円。
\item \textbf{音響管}：アクリル管にスピーカーで音を入れ、特定の位置に吸音材を置く。最も安い。コスト：\textbf{約1万円。1日で実験できる。}
\end{itemize}

\subsection{方法B：素数の篩を物質で作る（算術的境界条件）}

結晶の周期構造を使って、特定のモードだけを物理的にブロックする。

\begin{analogybox}
エラトステネスの篩を知っているだろうか？
2の倍数を消し、3の倍数を消し、5の倍数を消し…と繰り返すと、素数だけが残る。

これを「物質の結晶構造」として作る。
周期2の結晶で2の倍数をブロック、周期3で3の倍数をブロック…
重ねると「オイラー積結晶」＝物理的な篩が完成する。
\end{analogybox}

具体的な装置候補：
\begin{itemize}
\item \textbf{半導体超格子}：GaAs/AlGaAsを交互に積層。周期$p$の構造で$p$の倍数にバンドギャップが開く。
\item \textbf{メタマテリアル}：人工的な電磁材料。3方向に独立な周期を持たせれば、3つの素数を同時にミュート可能。
\end{itemize}

\subsection{方法C：ボスト-コンヌ系の量子シミュレーション}

最もダイレクトな方法。ボスト-コンヌ系のハミルトニアン
$H|n\rangle = \log(n)|n\rangle$ を量子コンピュータや冷却原子系で
そのまま再現し、素数の倍数の状態にペナルティを加える。

\begin{itemize}
\item \textbf{量子コンピュータ}：IBM QuantumやGoogle Cirqのクラウドで実行可能。\textbf{コスト：0円。1週間で実験できる。}
\item \textbf{冷却原子光格子}：レーザーで作った格子に超冷却原子を配置。基本格子に周期$p$の追加レーザーを重ねると、$p$の倍数のサイトだけ追加のポテンシャルを感じる。標準的なレーザー波長で実現可能。
\end{itemize}

\subsection{今すぐ始められる3つの実験}

\begin{keybox}[3つの即実行可能実験]
\begin{tabular}{lrrp{6cm}}
\toprule
実験 & コスト & 期間 & 何がわかるか \\
\midrule
$\alpha$：音響管 & 1万円 & 1日 & 算術的モード選別の原理が物理的に成立するか \\
$\beta$：量子コンピュータ & 0円 & 1週間 & BC系の素数ミュートが量子ハードウェアで再現できるか \\
$\gamma$：光学周波数コム & 100万円 & 6ヶ月 & 真空揺らぎの算術構造（オイラー積の物理的実在） \\
\bottomrule
\end{tabular}
\end{keybox}

大事なポイント：\textbf{$\alpha$と$\beta$は来週にでも始められる。}
物理学の最先端研究なのに、1万円と0円で始められる実験がある。
これは非常に珍しいことだ。

% ============================================================================
\section{壁の物理：ワープバブルの境界}
% ============================================================================

局所化を空間の一部だけに適用すると、
「WEC成立域」と「WEC不成立域」の境界ができる。
この境界がワープバブルの「壁」。

\begin{analogybox}
シャボン玉を想像してほしい。

シャボン玉の内側：$\Spec(\mathbb{Z}[1/2])$の世界。WEC不成立。ワープ可能。

シャボン玉の外側：$\Spec(\mathbb{Z})$の世界。WEC成立。通常の宇宙。

シャボン玉の膜：境界 = $\Spec(\mathbb{F}_2)$（2元体の1点スキーム）。
\end{analogybox}

この壁には面白い物理が予測される：

\begin{enumerate}
\item \textbf{エッジ状態}：
  トポロジカル絶縁体の表面に現れる「エッジ電流」のように、
  壁の上を走る特殊な量子状態が存在する。
  その数は $\mathbb{Q}_2/\mathbb{Z}_2$
  （2進分数全体）で数え上げられる無限個。

\item \textbf{放射}：
  ブラックホールの事象の地平面がホーキング放射を出すように、
  壁が微弱な熱放射を出す。
  温度は $\sim 30$ ミリケルビン（絶対零度の0.03度上）。

\item \textbf{屈折}：
  壁を通過する光が約30\%屈折する
  （壁の内外でモード密度が異なるため）。
  巨大な「重力レンズ」効果。
\end{enumerate}

% ============================================================================
\section{まとめ：何が問題で、何を解決したか}
% ============================================================================

\begin{table}[h]
\centering
\renewcommand{\arraystretch}{1.5}
\begin{tabular}{p{3cm}p{5cm}p{5cm}}
\toprule
& \textbf{旧い考え方}（カントール） & \textbf{新しい考え方}（グロタンディーク） \\
\midrule
真空とは & モードの「集合」 & $\Spec(\mathbb{Z})$ 上の「層」 \\
ミュートとは & 和から項を除去 & 空間の局所化 \\
WECとは & エネルギーの値 $\geq 0$ & レギュレータの「向き」 \\
符号反転とは & 値が $+$ から $-$ に変化 & 向きが「表」から「裏」に反転 \\
スケーリング問題 & $10^{69}$倍の増幅が必要 & \textbf{存在しない（予想）} \\
\bottomrule
\end{tabular}
\end{table}

\begin{keybox}[この論文の一言まとめ]
ワープドライブに必要な負のエネルギーは、
力ずくで「量」を集める問題ではなく、
数学の「構造」を変える問題かもしれない。

構造を変える道具が、グロタンディークの「層」と「局所化」。

鍵を壊すのではなく、鍵穴ごと消す。
\end{keybox}

% ============================================================================
\section{正直に言うと、まだ「予想」}
% ============================================================================

最後に、正直なことを言う。

この論文の核心的主張 ---
「WECはレギュレータの向きであり、スケーリング問題は消える」---
は\textbf{まだ証明されていない予想}だ。

証明されていないことの一覧：

\begin{enumerate}
\item 「WEC = レギュレータの向き」という対応関係自体が未証明。
\item 壁のエッジ状態や放射が本当に存在するかは未検証。
\item スケーリング問題が本当に消えるかは、予想が正しい場合に限る。
  もしWECが本質的に「エネルギーの量」の問題であれば、
  スケーリング問題は残る。
\end{enumerate}

しかし、数学において「未証明だが整合的な予想」は
研究の最も重要な原動力だ。
フェルマーの最終定理も、リーマン予想も、
長い間「予想」だった（前者は358年後に証明された）。

この論文が提示しているのは「答え」ではなく「新しい問いの立て方」。

WECは本当に「量」の条件なのか、それとも「構造」の条件なのか？

この問いに答えるために必要な実験が、
論文Aで提案されているSQUID実験
（既存技術で実施可能、約500万〜1000万円、12ヶ月）だ。

\vfill
\begin{center}
{\small Wright Brothers, 2026}
\end{center}

\end{document}
